
$X$ préschéma connexe, \struck{réduit, ce qui ne restreint pas}

1°) $X_{\mathrm{réd}}$ de car.\ $p > 0$. Alors dans le pro-torseur
fondamental $(T_\alpha)$, il n'y a \uncertain{pas} en général de sections
pour les $T_\alpha$.

On dispose uniquement, \struck{de \uncertain{façon} canonique}
\marginal{Coh.\ $\ell$-adique}
\begin{enumerate}
  \item[(i)] des $\xi_\ell^{(\alpha)} \in T_\alpha(\mathbb{Z}_\ell)$,
  $\ell \neq p$, provenant sans doute de $\xi_\ell \in
  T_\alpha(\widetilde{\mathbb{Z}_{\ell\mathbb{Z}}})$
  \marginal{hensélisé}
  \item[(ii)] d'un $\xi_p^{(\alpha)} \in T_\alpha(W_p(k))$ si
  \struck{$X$ déf.\ \ill{}} $k$ est un corps parfait sur $X$
  \struck{le corps parfait $k$}
  [$\xi_p$ provenant sans doute d'un $\xi_p \in
  T_p(\widetilde{\mathbb{Z}_{p\mathbb{Z}}})$]
  \marginal{hensélisé strict.}
  \item[(iii)] \uncertain{réduit} $\xi_p^{(\alpha)\,0} \in T_\alpha(k)$
  \marginal{« De Rham-Hodge »}
\end{enumerate}
\note{les items (ii) et (iii) sont réunis par une accolade, précédée d'un
« ? ». « hensélisé » et « hensélisé strict » sont écrits sous les deux
anneaux marqués d'un tilde, qu'ils glosent. « De Rham-Hodge » est écrit
au-dessus de (iii), avec une flèche vers le $k$ de $T_\alpha(k)$}

\marginal{Coh.\ $p$-adique. Structure supplémentaire : opérations $F$ et
$V$, ½ linéaires \uncertain{relativement} au Frobenius \ill{}}
\marginal{Coh.\ Hodge-de Rham. Mais ici, il faut \struck{sans doute}
travailler dans la catégorie \uncertain{dérivée}. Tenir compte aussi de la
\uncertain{filtration} Hodge-De Rham}
\note{ces notes marginales sont écrites en oblique dans la marge gauche, en
face de (i), (ii), (iii) respectivement : la deuxième est entourée d'un trait
qui la rattache à (ii) ; dans la deuxième, « ½ linéaires relativement au
Frobenius » est surchargé de plusieurs traits}

\uncertain{N.B.} Si $X = \operatorname{Spec} \mathbb{F}_q$, il y a en
général des \struck{\ill{}} obstructions non triviales dans des
$H^2(\mathbb{Q}, G)$ ---
\note{le $\mathbb{Q}$ de $H^2(\mathbb{Q}, G)$ est douteux ; le signe
initial, en marge, est une sorte de monogramme, lu N.B.}

On utilise des « foncteurs fibres » qui induisent sur les motifs effectifs
de poids $0$ le foncteur canonique. Alors les \uncertain{points} de
$T_\alpha(\mathbb{Z}_\ell)$ sont sans doute uniques [à isom.\ près] (th.\ de
Lang), de même ceux de $T_\alpha(W_p(k))$ si $k$ fini (ii), [plus précis.]
pour un $k$ parfait de dim.\ coh.\ $1$ (?). En général, il faut s'attendre
que $T_\alpha(k)$ \struck{\ill{}} ait plus d'un élément à isom.\ près, mais
il est possible que $T_\infty(k) = \varprojlim T_\alpha(k)$ n'ait qu'un
élément à isom.\ près, celui donné par la théorie de De Rham-Hodge.

\marginal{\uncertain{relation} \ill{} \uncertain{galoisienne} ?? Ce doit
être équivalent, cf.\ \emph{Manin}}
\note{note écrite en oblique dans la marge gauche, en bas, reliée par un
long trait de crayon à la note sur la cohomologie $p$-adique ; « Manin » est
souligné}

